\cleardoublepage
\frontmatterheading{Abstract}
\addcontentsline{toc}{chapter}{Abstract}

We experience indoor environments through our perceptions, emotions, prior experience, social context, and the wider conditions of everyday life. Human-Building Interaction (HBI) responds to this complexity by studying the relations between human experience, interactive technologies, and built environments. Yet the nature of this complexity remains under-specified: what kinds of experience matter indoors, how can HBI understand them, and how might this guide future interactive spaces? This thesis demonstrates how an \emph{affective lens} can address these questions by uncovering the nuances and forms of this complexity and orienting future design. It advances HBI through three connected contributions: affective interaction as a methodological lens to reveal lived experience in buildings, biophilic design as a way to translate this understanding into design approaches, and one AI-transformed indoor soundscape prototype as a concrete design instance.

The thesis begins with a scoping review of human experience design research in indoor environments, identifying emotional experience as a promising but under-examined way of approaching indoor complexity. It then investigates this potential through situated occupant studies, expert interviews, and expert design sessions. The occupant studies use Affective Interaction as a methodological lens to reveal how occupants sense, interpret, anticipate, and make meaning within socio-spatial contexts. The expert studies identify under-explored dimensions of biophilic experience in technologically mediated spaces and translate these possibilities into a design framework for HBI, with implications for HCI. Finally, an AI-transformed indoor soundscape prototype examines how everyday office sound can become nature-like and responsive, enriching quietness, reframing disruption, and opening possibilities for social, emotional, and biophilic experience indoors. Together, these studies position affective HBI as a way to understand, imagine, and design indoor environments from the complexity of human experience.